
\documentclass{article}
\usepackage{colortbl}
\usepackage{makecell}
\usepackage{multirow}
\usepackage{supertabular}

\begin{document}

\newcounter{utterance}

\centering \large Interaction Transcript for game `codenames', experiment `association\_difficult', episode 4 with reflexion{-}llama{-}3.1{-}8b{-}instruct.
\vspace{24pt}

{ \footnotesize  \setcounter{utterance}{1}
\setlength{\tabcolsep}{0pt}
\begin{supertabular}{c@{$\;$}|p{.15\linewidth}@{}p{.15\linewidth}p{.15\linewidth}p{.15\linewidth}p{.15\linewidth}p{.15\linewidth}}
    \# & $\;$A & \multicolumn{4}{c}{Game Master} & $\;\:$B\\
    \hline

    \theutterance \stepcounter{utterance}  
    & & & \multicolumn{2}{p{0.3\linewidth}}{
        \cellcolor[rgb]{0.9,0.9,0.9}{
            \makecell[{{p{\linewidth}}}]{
                \texttt{\tiny{[GM$|$GM]}}
                {'hidden': {'team': ['deacon', 'swan', 'gymnastics', 'france', 'nitrogen', 'hat'], 'innocent': ['ostrich', 'volleyball', 'korea', 'archbishop', 'acid'], 'opponent': [], 'assassin': ['coal']}, 'revealed': {'team': {'team': [], 'innocent': [], 'opponent': [], 'assassin': []}, 'opponent': {'team': [], 'innocent': [], 'opponent': [], 'assassin': []}}}
            }
        }
    }
    & & \\ \\

    \theutterance \stepcounter{utterance}  
    & & \multicolumn{4}{p{0.6\linewidth}}{
        \cellcolor[rgb]{0.9,0.9,0.9}{
            \makecell[{{p{\linewidth}}}]{
                \texttt{\tiny{[P1$\langle$GM]}}
                \texttt{You have attempted this same task before and did not succeed. These are the plans you wrote after those attempts.} \\
\texttt{Trial \#1: New plan: } \\
\\ 
\texttt{In my previous attempt, I provided a clue that was too specific and related to only two words. This led to the opponent guessing a word that was not a kill word but still ended the teammate's turn. Next time, I will provide a clue that relates to more words and is more general, but still semantically related to the target words. For example, if I notice that several words are related to a particular theme, I will choose a clue that is related to that theme. I will also make sure to check the lists of words that have been removed from the game to avoid providing a clue that is too specific to a word that is no longer in the game.} \\
\\ 
\texttt{Trial \#2: New plan: } \\
\\ 
\texttt{In my previous attempt, I provided a clue that was too specific and related to only two words, but still managed to fail to guess all the words before the opponent. This time, I will provide a clue that relates to more words and is more general, but still semantically related to the target words. However, I also need to pay more attention to the words that have been removed from the game and avoid providing a clue that is too specific to a word that is no longer in the game. } \\
\\ 
\texttt{In my last attempt, I provided a clue that was related to the theme of "fashion" with the clue "fabric" and then "tree" which is a more general theme. However, I failed to guess all the words before the opponent. This time, I will try to find a clue that relates to a more general theme that encompasses more words. For example, if I notice that several words are related to a particular theme, such as "nature" or "animals", I will choose a clue that is related to that theme. I will also make sure to check the lists of words that have been removed from the game to avoid providing a clue that is too specific to a word that is no longer in the game.} \\
\\ 
\\ 
\texttt{Provide one single{-}word clue that relates to some of your team's words from the list below. You can choose to how many and to which words your clue relates to.} \\
\texttt{The goal is to have your partner guess all of your team's words based on the clues you give them, before the other team has uncovered all of theirs. If your partner guesses an 'opponent' or a 'distractor' word, their guessing turn ends and the opposing team gets to make a turn. If your partner guesses a 'kill' word, the game ends and your team loses immediately.} \\
\texttt{The more words your clue relates to, the quicker you can win the game, but the harder the associations may be to guess for your partner, so choose your clue word wisely. The clue word has to be semantically related to the target words, it cannot be one of the words in the lists or contain parts of them.} \\
\\ 
\texttt{Always give your single{-}word clue and your comma{-}separated list of related target words in the following format and make your answers as short as possible, never include any other text than is required in this form:} \\
\\ 
\texttt{CLUE: <WORD>} \\
\texttt{TARGETS: <TARGETS>} \\
\\ 
\texttt{Your team words are: deacon, swan, gymnastics, france, nitrogen, hat.} \\
\texttt{Your opponent's team words are: .} \\
\texttt{Distractor words are: ostrich, volleyball, korea, archbishop, acid.} \\
\texttt{The kill words are: coal.} \\
            }
        }
    }
    & \\ \\

    \theutterance \stepcounter{utterance}  
    & \multicolumn{4}{p{0.6\linewidth}}{
        \cellcolor[rgb]{0.9,0.9,0.9}{
            \makecell[{{p{\linewidth}}}]{
                \texttt{\tiny{[P1$\rangle$GM]}}
                \texttt{CLUE: church} \\
\texttt{TARGETS: deacon, france} \\
            }
        }
    }
    & & \\ \\

    \theutterance \stepcounter{utterance}  
    & & & \multicolumn{2}{p{0.3\linewidth}}{
        \cellcolor[rgb]{0.9,0.9,0.9}{
            \makecell[{{p{\linewidth}}}]{
                \texttt{\tiny{[GM$|$GM]}}
                \texttt{church} \\
            }
        }
    }
    & & \\ \\

    \theutterance \stepcounter{utterance}  
    & & & \multicolumn{2}{p{0.3\linewidth}}{
        \cellcolor[rgb]{0.9,0.9,0.9}{
            \makecell[{{p{\linewidth}}}]{
                \texttt{\tiny{[GM$|$GM]}}
                ['deacon', 'france']
            }
        }
    }
    & & \\ \\

    \theutterance \stepcounter{utterance}  
    & & & \multicolumn{2}{p{0.3\linewidth}}{
        \cellcolor[rgb]{0.9,0.9,0.9}{
            \makecell[{{p{\linewidth}}}]{
                \texttt{\tiny{[GM$|$GM]}}
                {'word': 'deacon', 'assignment': 'team'}
            }
        }
    }
    & & \\ \\

    \theutterance \stepcounter{utterance}  
    & & & \multicolumn{2}{p{0.3\linewidth}}{
        \cellcolor[rgb]{0.9,0.9,0.9}{
            \makecell[{{p{\linewidth}}}]{
                \texttt{\tiny{[GM$|$GM]}}
                {'word': 'france', 'assignment': 'team'}
            }
        }
    }
    & & \\ \\

    \theutterance \stepcounter{utterance}  
    & & & \multicolumn{2}{p{0.3\linewidth}}{
        \cellcolor[rgb]{0.9,0.9,0.9}{
            \makecell[{{p{\linewidth}}}]{
                \texttt{\tiny{[GM$|$GM]}}
                \texttt{opponent has won} \\
            }
        }
    }
    & & \\ \\

    \theutterance \stepcounter{utterance}  
    & \multicolumn{4}{p{0.6\linewidth}}{
        \cellcolor[rgb]{0.95,0.95,0.95}{
            \makecell[{{p{\linewidth}}}]{
                \texttt{\tiny{[P2$\rightarrow$P1]}}
                \texttt{Post{-}episode Reflexion (not sent to GM):} \\
\texttt{New plan: } \\
\\ 
\texttt{In my previous attempts, I provided clues that were too specific and related to only a few words, which led to the opponent guessing a word that ended the teammate's turn. This time, I will provide a clue that relates to a more general theme that encompasses more words. I will also pay closer attention to the words that have been removed from the game and avoid providing a clue that is too specific to a word that is no longer in the game. } \\
\\ 
\texttt{Upon reviewing the list of words, I notice that several words are related to a theme of "religion" or "spirituality", such as "deacon" and "church". I will choose a clue that is related to this theme, but is general enough to encompass more words. For example, I will choose a clue that is related to a concept or idea that is common to several words on the list.} \\
            }
        }
    }
    & & \\ \\

\end{supertabular}
}

\end{document}
